\begin{lemma}
Let $V,T$ be finite sets, $s:V\to T$ injective, and $0\le p\le1$.
Put $b_p=p\delta_1+(1-p)\delta_0$, $u=\operatorname{Leb}|_{[0,1]}$,
and $m=(b_p\otimes u)^T$. Let $\mu$ be the push-forward of $m$ by
$q\mapsto(\{v:q(s(v))_1=1\},(q(s(v))_2)_v)$.
Then for every $A\subseteq V$,
\[
\mu\{(B,\pi):B=A\}=\operatorname{ofReal}(p^{|A|}(1-p)^{|V|-|A|}),
\]
and for every $t:V\to[0,1]$,
\[
\mu\{(B,\pi):B=A,\ 0\le\pi_v\le t_v\ (v\in V)\}
=\mu\{B=A\}\operatorname{ofReal}(\prod_v t_v).
\]
\end{lemma}
\begin{proof}
The decoding map is measurable by \noderef{SamplingCoordinateDecodeMeasurable}.
Its displayed events are measurable, since their activation fibers are discrete
and their priority conditions are finite intersections of closed coordinate
conditions. Thus their push-forward masses equal the masses of their inverse
images. The measure $b_p$ has total mass $p+(1-p)=1$, with mass $p$ on
$\{1\}$ and $1-p$ on its complement, because $0\ne1$.
Also $u(\mathbb R)=1$ and $u([0,a])=a$ for $0\le a\le1$.
For the activation event, its inverse image is the product rectangle whose
factor at $s(v)$ is $\{1\}\times\mathbb R$ if $v\in A$, and
$(\mathbb R\setminus\{1\})\times\mathbb R$ otherwise; at tags not in
$s(V)$ its factor is $\mathbb R^2$. Injectivity makes these assignments
unambiguous. Finite-product rectangle evaluation therefore gives one factor
$p$ for each member of $A$, one factor $1-p$ for each member of $V\setminus A$,
and a factor one for every unused tag. This is the first formula.
For the second event replace each used tag's priority factor $\mathbb R$
by $[0,t_v]$. Product-measure rectangle evaluation now multiplies the previous
activation factors by $t_v$ for each $v$. Regrouping the finite product yields
the first mass times $\prod_v t_v$. All factors are nonnegative, so passage
through $\operatorname{ofReal}$ respects these products. No division was used;
the argument includes $p=0,1$ and empty coordinate sets, with empty products one.
\end{proof}
